% LTeX: enabled=false

\documentclass{article}
\usepackage[utf8]{inputenc}

\newcommand{\nirpdftitle}{Lie Theory Review}
\usepackage{import}
\inputfrom{..}{nir}

\pagestyle{contentpage}

\title{Lie Theory for the Impatient}
\author{Nir Elber}
\date{Summer 2026}
\rhead{\textit{LIE THEORY SPEEDRUN}}
\lhead{}

\setcounter{tocdepth}{2}

\begin{document}

\maketitle

\begin{abstract}
	This document collects a variety of definitions and results from Lie theory.\todo{quaternionic}
\end{abstract}

\tableofcontents

\newpage

\section{Topics in Representation Theory}

\subsection{Miniscule Weights}
\begin{definition}[miniscule]
	A dominant integral weight $\omega$ is \textit{miniscule} if and only if $(\omega,\beta)\le1$ for all positive coroots $\beta$.
\end{definition}
\begin{remark}
	Nonzero miniscule weights are fundamental: $(\omega,\alpha_i^\lor)\le1$, with equality for at most one $i$.
\end{remark}
\begin{proposition}
	Let $\omega$ be a dominant integral weight. The following are equivalent.
	\begin{listroman}
		\item $\omega$ is miniscule.
		\item All weights of $L_\omega$ belong to $W\omega$.
		\item If $\lambda$ is dominant integral, and $\omega-\lambda\in Q^+$, then $\lambda=\omega$.
	\end{listroman}
\end{proposition}
\begin{proof}[Sketch]
	For (iii) implies (ii), note a weight $\lambda$ of $L_\omega$ has $w\lambda$ dominant for some $w$, but then (iii) means $w\lambda=\omega$. For (ii) implies (i), suppose (i) is false, and we hunt for a suitable weight space in $L_\omega$. Well, choose a root $\alpha$ with $(\omega,\alpha)>1$, and check that $\omega-\alpha$ is a weight space of $L_\omega$ not in $W\omega$.

	It remains for (i) implies (iii), which is some difficult combinatorics that I do not understand.\todo{}
\end{proof}
\begin{proposition}
	Suppose $\omega$ is miniscule, and $\lambda$ is dominant integral. Then
	\[L_\omega\otimes L_\lambda=\bigoplus_{\substack{\gamma\in W\omega\\\lambda+\gamma\ge0}}L_{\lambda+\gamma}.\]
\end{proposition}
\begin{proof}[Sketch]
	Use the Weyl character formula to expand both sides.
\end{proof}
\begin{proposition}
	Each coset in $P/Q$ admits a unique miniscule weight.
\end{proposition}
\begin{proof}
	In one direction, suppose $\omega$ is the dominant integral weight in a coset of $P/Q$ of minimal height. Then the only allowed weights of $L_\omega$ are in $W\omega$ for height reasons, so $\omega$ miniscule.

	In the other direction, if $\omega$ and $\omega'$ are miniscule weights in the same coset of $P/Q$, we know $\omega_1-\omega_2\in Q$. To show $\omega_1=\omega_2$, note that $(\omega_1-\omega_2,\beta)\le1$ for each coroot $\beta$: we may assume $\beta$ is positive, and then $(\omega_1,\beta)$ and $(\omega_2,\beta)$ are both in $\{0,1\}$. But if $\omega\in Q$ has $(\omega,\beta)\le1$ for all coroots $\beta$, we must have $\omega=0$: write $\omega=\sum_im_i\alpha_i$ as the minimal counterexample, but one can show that $s_j\omega=\omega-\alpha_j$ is a smaller counterexample for some $j$ (choose $j$ so that $m_j$ and $(\omega,\alpha_j^\lor)$ have the same sign, which we may assume is positive).
\end{proof}

\subsection{Representations of \texorpdfstring{$\op{GL}_n$}{GLn}}
\begin{theorem}[Schur--Weyl duality]
	Let $V$ be a complex vector space of dimension $n$. Then there is a decomposition
	\[V^{\otimes N}=\bigoplus_{\substack{\lambda=(\lambda_1,\ldots,\lambda_n)\\\lambda\vdash N}}L_\lambda\otimes\pi_\lambda\]
	into irreducible representations of $\op{GL}_n(\CC)\times S_N$, where $L_\lambda$ is the irreducible representation of $\op{GL}(V)$ of highest weight $\lambda$. If $n\ge N$, the the set $\{\pi_\lambda\}_\lambda$ exhausts the set
\end{theorem}
\begin{proof}[Sketch]
	One uses the double-centralizer lemma on the following subalgebras of $\op{End}_\CC V^{\otimes N}$: define $A$ to be the image of $U\mf{gl}_n(\CC)$, and define $B$ to be the image of $\CC[S_N]$. For the last assertion, one notes that the quotient map $\CC[S_N]\to B$ is an isomorphism.
\end{proof}
\begin{definition}
	For a partition $\lambda\vdash N$ and a vector space $V$ over $\CC$, define the \textit{Schur functor}
	\[S^\lambda V\coloneqq\op{Hom}_{S_N}\left(\pi_\lambda,V^{\otimes N}\right).\]
\end{definition}
\begin{theorem}[Invariant theory]
	Fix a finite-dimensional vector space $V$. Then we can describe a spanning set of $\op{GL}(V)$-invariant functions on the space
	\[\bigotimes_i\left(V^{\otimes m_i}\otimes V^{*\otimes n_i}\right)^{\otimes d_i}\]
	as follows: start with $d_i$ vertices for each $i$, required to have $m_i$ incoming and $n_i$ outgoing arrows. For each graph $\Gamma$ built from such vertices, define $F_\Gamma$ as the associated convolution.\todo{}
\end{theorem}
\begin{proof}[Sketch]
	We are hunting for $\op{GL}(V)$-invariant functionals on the given tensor power. By rearranging, this should amount to calculating $\op{GL}(V)$-invariant endomorphisms of $V^{\otimes N}$ for some $N$, which are spanned by permutations by Schur--Weyl duality.
\end{proof}
\begin{definition}
	Fix a partition $\lambda$ with at most $n$ parts. By the Weyl character formula, the character of the irreducible representation $L_\lambda$ of $\op{GL}_n(\CC)$ is given by the \textit{Schur polynomial}
	\[s_\lambda(x_1,\ldots,x_n)=\frac{\det\left[x_i^{\lambda_j+n-j}\right]_{ij}}{\displaystyle\prod_{i<j}(x_i-x_j)}.\]
\end{definition}
\begin{theorem}[Frobenius character formula]
	Suppose $\sigma\in S_N$ is of cycle type $(\ell_1,\ldots,\ell_r)$. Then for a partition $\lambda$ of $N$, the character value $\tr(\sigma;\pi_\lambda)$ is the coefficient of $x_1^{\lambda_1+-n-1}\cdots x_n^{\lambda_n}$ of the polynomial
	\[\prod_{i<j}(x_i-x_j)\cdot\prod_{i=1}^r\left(x_1^{\ell_i}+\cdots+x_n^{\ell_r}\right).\]
\end{theorem}
\begin{proof}[Sketch]
	For diagonal $x\in\op{GL}_n(\CC)$, one can compute the trace of $x\sigma$ acting on $V^{\otimes N}$ in two ways: direct calculation on the standard basis of $V^{\otimes N}$ shows that
	\[\tr\left(x\sigma;V^{\otimes N}\right)=\prod_{i=1}^r\left(x_1^{\ell_i}+\cdots+x_n^{\ell_i}\right).\]
	Alternatively, we can compute the trace of $x\sigma$ in the decomposition of $V^{\otimes N}$ given by Schur--Weyl duality.
\end{proof}
\begin{theorem}[Howe duality]
	For vector spaces $V$ and $W$, we have
	\[S^N(V\otimes W)=\bigoplus_{\lambda\vdash N}S^\lambda V\otimes S^\lambda W.\]
\end{theorem}
\begin{proof}[Sketch]
	The key idea is that $S^N(V\otimes W)=\left(V^{\otimes N}\otimes W^{\otimes N}\right)^{S_N}$. One now expands $S^{\otimes N}$ and $W^{\otimes N}$ using Schur--Weyl duality. Eventually, one needs to know that $(\pi_\lambda\otimes\pi_\mu)^{S_N}$ is at most one-dimensional (by Schur's lemma), and it is one-dimensional if and only if $\lambda=\mu$ (because $\pi_\mu$ is self-dual because its character is real-valued).
\end{proof}
\begin{proposition}
	Let $\lambda$ be a partition of $N$. Then
	\[\CC[S_{N+1}]\otimes_{\CC[S_N]}\pi_\lambda=\bigoplus_{\mu\in\lambda+\square}\pi_\mu.\]
\end{proposition}
\begin{proof}[Sketch]
	It is enough to check that both representations give the same multiplicity space in $V^{\otimes(N+1)}$ for $V$ of large dimension. The right-hand side gives $V\otimes S^\lambda V$ because $V$ is miniscule of weight $(1,0,\ldots,0)$. The left-hand side gives $V\otimes S^\lambda V$ by applying Frobenius reciprocity and then Schur--Weyl duality.
\end{proof}
\begin{corollary}
	Let $\lambda$ be a partition. Then $\pi_\lambda\otimes\mathrm{sgn}=\pi_{\lambda^\dagger}$.
\end{corollary}
\begin{proof}[Sketch]
	Induct on $\left|\lambda\right|$. Then the proposition implies
	\[\bigoplus_{\mu\in\lambda+\square}\pi_\mu\otimes\mathrm{sgn}=\bigoplus_{\mu\in\lambda^\dagger+\square}\pi_\mu.\]
	One now uses the Casimir element of $\op{GL}_n(\CC)$ (equivalently, the Jucys--Murphy element of $\CC[S_{N+1}]$ given by $\sum_{i<j}(i,j)$) to separate out the representations appropriately by eigenvalue.
\end{proof}

\subsection{The Automorphism Group}
\begin{definition}[adjoint group]
	For a finite-dimensional Lie algebra $\mf g$, let $G_{\mathrm{ad}}$ be $\op{Aut}(\mf g)^\circ$.
\end{definition}
\begin{remark}
	A computation on differentials shows $\op{Lie}G_{\mathrm{ad}}=\op{Der}\mf g$. If $\mf g$ is semisimple, then the natural map $\mf g\to\op{Der}\mf g$ is an isomorphism, so $\op{Lie}G_{\mathrm{ad}}=\mf g$.
\end{remark}
\begin{lemma}
	Let $\mf h$ be a Cartan subalgebra, and let $T\subseteq G_{\mathrm{ad}}$ be the corresponding subgroup. If $h\in\op{Aut}(\mf g)$ fixes $\mf h$ and the root spaces, then $h\in T$.
\end{lemma}
\begin{proof}[Sketch]
	For each simple root $\alpha_i$, find $c_i\in\CC$ such that the action of $h$ on $\mf g[\alpha_i]$ is $\exp(c_i)$. Then the automorphism $h$ of $\mf g$ equals the automorphism $\exp\left(\sum_ic_i\omega_i^\lor\right)$.
\end{proof}
\begin{lemma}
	Let $\mf h$ be a Cartan subalgebra, and let $T\subseteq G_{\mathrm{ad}}$ be the corresponding subgroup. If $g\in G_{\mathrm{ad}}$ preserves $\mf h$ and a set of simple roots, then $g=1$.
\end{lemma}
\begin{proof}[Sketch]
	The collection of $g\in G_{\mathrm{ad}}$ which preserve $\mf h$ and a set of simple roots must then permute the simple root spaces and hence induces an automorphism $\sigma_g$ of the Dynkin diagram. The idea is that this automorphism of the Dynkin diagram should be reflected in the fundamental representations, which are discrete, so $g=1$ will be forced by $G_{\mathrm{ad}}$ being connected.

	In a few more words, note $g\in G_{\mathrm{ad}}$ acts on $\op{Rep}\mf g$ by sending $\rho$ to $\rho\circ\op{Ad}_{g^{-1}}$.
	\begin{itemize}
		\item On one hand, if $g$ preserves $\mf h$ and a set of simple roots (so that it preserves $\mf n_+$), we see that $g$ preserves singular vectors, so for any dominant integral $\lambda$, we see $\rho_\lambda\circ\mathrm{Ad}_{g^{-1}}=\rho_{\sigma\lambda}$.
		\item On the other hand, this action is essentially trivial: by the Weyl character formula, for any dominant integral weight $\lambda$, we see $g\mapsto\tr(-;\rho_\lambda\circ\mathrm{Ad}_{g^{-1}})$ is a continuous family of polynomials valued in the characters. But the characters are a discrete set of polynomials, so this continuous family is constant. We conclude $\rho_\lambda\circ\mathrm{Ad}_{g^{-1}}=\rho_\lambda$.
		\qedhere
	\end{itemize}
\end{proof}
\begin{proposition}
	Let $\mf h$ be a Cartan subalgebra, and let $T\subseteq G_{\mathrm{ad}}$ be the corresponding subgroup. Then $N(T)/T\cong W$.
\end{proposition}
\begin{proof}[Sketch]
	The main point is to construct lifts of the simple reflections of $W$. Accordingly, for each simple root $\alpha_i$, let $\dot s_i\in G_{\mathrm{ad}}$ be the image of $\begin{bsmallmatrix}
		0 & 1 \\ -1 & 0
	\end{bsmallmatrix}$ along the map $\op{SL}_2(\CC)\to G_{\mathrm{ad}}$ defined by lifting $\mf{sl}_{2i}\subseteq\mf g$. Then $\op{Ad}_{\dot s_i}$ sends $\alpha_i\mapsto-\alpha_i$ and fixes $\{\alpha_i\}^\perp$ (because the full one-parameter subgroup $t\mapsto\exp(th_i)$ does), so the action agrees with the action of $s_i$ on $\mf h$. Thus, $\dot s_i$ stabilizes $\mf h$ and lives in $N(T)$.

	Extending this construction multiplicatively produces a subgroup $N\subseteq N(T)$ such that $N/T\cong W$. To see $N=N(T)$, note that any $g\in N(T)$ preserves $\mf h$, and we may adjust $g$ by $N\onto W$ so that $g$ preserves a set of simple roots of $\mf h$, which forces $g=1$.
\end{proof}
\begin{proposition}
	Let $D$ by the Dynkin diagram, acting on $\mf g$ by permuting the Serre presentation. Then the natural map
	\[\op{Aut}D\to\pi_0\op{Aut}(\mf g)\]
	is an isomorphism.
\end{proposition}
\begin{proof}[Sketch]
	We already know $G_{\mathrm{ad}}\cap\op{Aut}D$ is trivial, so the difficulty is surjectivity. Well, for $g\in\op{Aut}(\mf g)$, choose a $g$-stable Cartan subalgebra $\mf h\subseteq\mf g$. Adjusting by $W$, we may assume $g$ preserves a set of simple roots. Adjusting by $\op{Aut}D$, we may assume $g$ preserves the set of simple roots pointwise. Adjusting by $H$, we may assume $g$ preserves the root spaces pointwise. It follows $g=1$. 
\end{proof}

\subsection{Miscellaneous Topics}
\begin{proposition}
	Let $w_0$ be the long Weyl element. For any dominant integral weight $\lambda$, we have $L_\lambda^*=L_{-w_0\lambda}$.
\end{proposition}
\begin{proof}[Sketch]
	The lowest weight of $L_\lambda$ is $w_0\lambda$, which we see by considering the Weyl orbit. Negate to get a dominant weight.
\end{proof}
\begin{definition}[maximal root]
	The \textit{maximal root} $\theta$ is the dominant integral weight associated to the adjoint representation, which is a root by the root decomposition.
\end{definition}
\begin{remark}
	Because $L_\theta$ is a quotient of the Verma module $M_\theta$, we see that any root $\alpha$ has $\theta-\alpha\in Q^+$. It follows that the height of $\theta$ is maximal among roots.
\end{remark}
\begin{remark}
	The maximal root is always long, easily seen from the classification.
\end{remark}
\begin{remark}
	The maximal root is a fundamental weight outside of types $A$ and $C$, which is checked from the classificaiton.
\end{remark}
\begin{definition}[Coxeter number]
	The \textit{Coxeter number} is $(\theta,\rho^\lor)+1$. This is one more than the height of the maximal root. The \textit{dual Coxeter number} is $(\widetilde\theta^\lor,\rho)+1$, where $\widetilde\theta^\lor$ is the coroot of $\theta$.
\end{definition}
\begin{remark}
	If $\mf g$ is simply laced, the dual root system agrees withthe root system, so $\widetilde\theta^\lor=\theta^\lor$, so the Coxeter number and dual Coxeter number are equal. This is false without the simply laced assumption.
\end{remark}
\begin{definition}[principle $\mf{sl}_2$-triple]
	The \textit{principal $\mf{sl}_2$-triple} $(e,f,h)$ has $e=\sum_ie_i$ and $h=2\rho^\lor$ and $f=\sum_i(2\rho^\lor,\omega_i)f_i$.
\end{definition}
\begin{definition}[exponents]
	Consider $\mf g$ as a representation of the principle $\mf{sl}_2$. For $m\ge0$, define $r_m\coloneqq\dim\mf g[2m]$, which is the number of (positive) roots of height $m$ if $m>0$. Then the \textit{exponents} are the positive integers $m$ for which $r_m>r_{m+1}$, and the multiplicity is $r_m-r_{m+1}$.
\end{definition}
\begin{example}
	We have $r_0=\op{rank}\mf g$ by the root decomposition, $r_1=\op{rank}\mf g$ is the number of simple roots, and $r_2=\op{rank}\mf g-1$ because $\alpha_i+\alpha_j$ is a root if and only if $\alpha_i$ and $\alpha_j$ are adjacent in the Dynkin diagram (seen from the classification of rank-$2$ root systems).
\end{example}
\begin{definition}[type]
	Let $V$ be a finite-dimensional representation. Then $V$ is
	\begin{itemize}
		\item of \textit{complex type} if and only if $V\not\cong V^*$,
		\item of \textit{real type} if and only if $V\cong V^*$ by a symmetric bilinear form, and
		\item of \textit{quaternionic type} if and only if $V\cong V^*$ by an alternating bilinear form.
	\end{itemize}
\end{definition}
\begin{proposition}
	Let $\lambda$ be a dominant integral weight. Then
	\begin{listalph}
		\item is of complex type if and only if $\lambda\ne-w_0\lambda$,
		\item is of real type if and only if $\lambda=-w_0\lambda$ and $(2\rho^\lor,\lambda)$ is even, and
		\item of quaternioic type if and only if $\lambda=-w_0\lambda$ and $(2\rho^\lor,\lambda)$ is odd.
	\end{listalph}
\end{proposition}
\begin{proof}[Sketch]
	Here, (a) follows because $L_\lambda^*=L_{-w_0\lambda}$. One checks (b) and (c) by passing to the principal $\mf{sl}_2$-triple, which features a highest weight vector $v_\lambda$ of multiplicity one by height considerations.
\end{proof}

\section{The Peter--Weyl Theorem}

\subsection{Measure Theory}
\begin{theorem}[Haar]
	Let $G$ be a locally compact group. Then $G$ admits a unique (up to scalar) left-invariant Radon measure $\mu$. There is an analogous statement for right-invariance.
\end{theorem}
\begin{proof}[Sketch]
	% Uniqueness is easier: given two such measures $\mu$ and $\nu$, let $I_\mu$ and $I_\nu$ be the associated positive continuous linear functionals on $C_c(G;\RR)$. It is enough to show that
	% \[I_\mu(f)I_\nu(g)=I_\mu(g)I_\nu(f)\]
	% for any $f,g\in C_c(G;\RR)$. For an approximate identity $\{\varphi_n\}$ of $G$, note that $\varphi_n\cdot f\to f$, so
	% \[\int_G\varphi_n(x)f(xy)\,d\mu(x)=f(x)+o(1).\]
	% Thus,
	% \[\int_G\int_G\varphi_n(x)f(xy)\,d\mu(x)\,d\nu(y)=I_\mu(f)\]
	We ignore uniqueness. Roughly speaking, it follows from the Fubini--Tonelli theorem, which we see by comparing double integrals over $\mu$ and $\nu$.\todo{}

	The difficulty lies in existence. We only discuss some special cases.
	\begin{itemize}
		\item If $G$ is discrete, then the counting measure suffices.
		\item Suppose $G$ is a Lie group (or more generally the topological space associated to the rational points of an algebraic group over a local field). If $G$ is a Lie group, one can construct measures by choosing a global section $\xi$ of $\det TG$. To have left-invariance, we must have $\xi_g=L_g^*\xi_1$, so it is enough to choose $\xi_1\in\det T_eG$.
		\item Suppose $G$ is compact and separable. Let $\{g_i\}$ be a dense subsequence, and choose positive weights $\{p_i\}$ so that $\sum_ip_i=1$. It is enough to exhibit an invariant positive continuous functional $I$ on $C(G;\RR)$. The idea is to view $I$ as the ``average'' of a function, which we can compute iteratively with the approximation
		\[Af(x)\coloneqq\sum_ip_if(xg_i).\]
		In particular, $A$ is a bounded linear operator of norm one; also, setting $v(f)\coloneqq\max f-\min f$ to be the variation, if $v(f)>0$, one has that $\max Af<\max f$ and $\min Af>\min f$ and so $v(Af)<v(f)$. One can use the Arzel\`a--Ascoli theorem to show that $\{A^nf\}_n$ admits a convergent subsequence, and we see that the limit must be constant, which one takes to be $I(f)$. (In fact, $\{A^nf\}_n$ converges: the sequences $\{\max A^nf\}_n$ and $\{\min A^nf\}_n$ are monotone, bounded, and they have a subsequence converging to $I(f)$.) % It is now straightforward to check that $I$ has the required properties; for example, $I$ is left-invariant because $A$ is.
		\qedhere
	\end{itemize}
\end{proof}
\begin{definition}[unimodular]
	A locally compact group $G$ is \textit{unimodular} if and only if left-invariant Radon measure on $G$ are also right-invariant.
\end{definition}
\begin{proposition}
	Fix a (real) Lie group $G$. Then $G$ is unimodular if and only if the character $\left|\det\mathrm{Ad}\right|$ is trivial.
\end{proposition}
\begin{proof}[Sketch]
	Choose nonzero $\xi\in\det\mf g$. We want to know if the left-invariant global section $g\mapsto\left|L_g^*\xi\right|$ is also right-invariant. It is equivalent to ask if $g\mapsto\left|R_{g^{-1}}^*L_g\xi\right|$ is trivial, which is equivalent to ask if $\left|\xi\right|$ is conjugation-invariant.
\end{proof}
\begin{example}
	If $G$ is compact, then $G$ is unimodular: any continuous character $G\to\RR^+$ has compact and hence trivial iamge.
\end{example}
\begin{definition}
	A Lie algebra $\mf g$ is unimodular if and only if $\det\op{ad}$ is the trivial representattion.
\end{definition}
\begin{remark}
	If $G$ is connected, then the character $\det\op{Ad}\colon G\to\RR^\times$ has connected image and hence already outputs to $\RR^+$. Thus, $G$ is unimodular if and only if $\mf g$ is unimodular.
\end{remark}
\begin{example}
	If $\mf g$ is abelian, it is unimodular because $\op{ad}$ is zero. If $\mf g$ is semisimple, then $[\mf g,\mf g]=\mf g$, so $\mf g$ is unimodular because any character $\mf g\to\CC$ is zero. Taking the sum shows that $\mf g$ is unimodular if $\mf g$ is reductive.
\end{example}

\subsection{Matrix Coefficients}
\begin{proposition}
	Let $G$ be a compact group. Then any finite-dimensional representation $V$ of $G$ is unitary.
\end{proposition}
\begin{proof}[Sketch]
	Choose any Hermitian form on $V$, and then average it over $G$ using integration to produce a $G$-invariant Hermitian form.
\end{proof}
\begin{proposition}
	Let $G$ be a Lie group. For any continuous finite-dimensional representation $V$, the matrix coefficients of $V$ are smooth.
\end{proposition}
\begin{proof}[Sketch]
	Let $V^{\mathrm{sm}}$ be the subspace of $V$ consisting of $v\in V$ for which any matrix coefficient $g\mapsto\ell(gv)$ is smooth. Because convolution picks up the best regularity properties, we know that $\varphi\cdot v$ is smooth for any $v\in V$ and $\varphi\in C^\infty(G)$. By using an approximate identity, one finds that $V^{\mathrm{sm}}$ is dense in $V$. Thus, $V^{\mathrm{sm}}=V$ because $V$ is finite-dimensional.
\end{proof}
\begin{proposition}[Orthogonality]
	Let $G$ be a compact group, and let $V$ and $W$ be non-isomorphic irreducible finite-dimen\-sional representations. For $v\in V$ and $v^*\in V^*$ and $w\in W$ and $w^*\in W^*$, we have
	\[\int_G\langle gv,v^*\rangle\overline{\langle gw,w^*\rangle}\,dg=0.\]
\end{proposition}
\begin{proof}[Sketch]
	Because $W$ is unitary, $\overline W=W^*$. Thus, it is enough to show
	\[\int_G\langle gv,v^*\rangle\langle w,gw^*\rangle\,dg=0.\]
	under the assumption $V\not\cong W^*$. Well, fix $v^*$ and $w$, and then the associated map $V\otimes W\to\CC$ is $G$-invariant and hence vanishing unless $V\cong W^*$.
\end{proof}
\begin{proposition}[Orthogonality]
	Let $G$ be a compact group, and let $V$ be an irreducible finite-dimensional representation. For $v,w\in V$ and $v^*,w^*\in V^*$, we have
	\[\int_G\langle gv,v^*\rangle\overline{\langle gw,w^*\rangle}\,dg=\frac1{\dim V}\langle v,w^*\rangle\langle w,v^*\rangle.\]
\end{proposition}
\begin{proof}[Sketch]
	Upon identifying $\overline V$ with $V^*$, the left-hand side is
	\[\int_G\langle gv,v^*\rangle\langle w,gw^*\rangle\,dg.\]
	If one fixes $v^*$ and $w$, this is an invariant Hermitian form in $v$ and $w^*$, which is unique up to scalar, which explains the factor of $\langle v,w^*\rangle$ on the right. The remaining scalar is still a Hermitian form in $(v,w^*)$, which is unique up to scalar again and explais the factor of $\langle w,v^*\rangle$ on the right. To get the remaining scalar $1/\dim V$, sum both sides over an orthnormal basis of $V$.
\end{proof}
\begin{definition}[finite]
	For a continuous representation $V$ of a compact group $G$, let $V^{\mathrm{fin}}$ be the subspace of $G$-finite vectors, where $v$ is \textit{$G$-finite} if and only if $\op{span}\{gv:g\in G\}$ is finite-dimensional.
\end{definition}
\begin{corollary}
	Let $G$ be a compact group. Consider the matrix coefficient map
	\[\xi\colon\bigoplus_V(V\otimes V^*)\to L^2(G)\]
	given on pure tensors by sending $v\otimes v^*$ to $g\mapsto\langle v,gv^*\rangle$. Then $\xi$ is $(G\times G)$-invariant, an isometric embedding, and has image given by $L^2(G)^{\mathrm{fin}}$.
\end{corollary}
\begin{proof}[Sketch]
	Here, $G\times G$ acts on $C(G)$ by left and right translation, so invariance is checked directly. The isometry property follows from the orthogonality relations.
	
	As for the image, certainly the image contains only $G$-finite vectors. Conversely, if $\varphi\in L^2(G)$ is $G$-finite, we may as well suppose that $\varphi$ generates an irreducible representation $V$, so $\varphi$ lives in a multiplicity space of $\op{Hom}_G\left(V,L^2(G)\right)$. Thus, it is enough to show that the matrix coefficient map $V^*\to\op{Hom}_G\left(V,L^2(G)\right)$ is an isomorphism: it turns out that any $G$-invariant map $V\to L^2(G)$ is uniquely determined by one value, which can then be spread out to a canonical invariant map $V\to C(G)$, and then there is an inverse map sending sending $\varphi\colon V\to L^2(G)$ to the functional $v\mapsto\varphi(v)(1)$.\todo{rewrite to show $K$-finite implies smooth}
\end{proof}

\subsection{Functional Analysis}
\begin{definition}[compact]
	A bounded linear operator $A$ on a Hilbert space $H$ is \textit{compact} if and only if it is the limit (in the norm topology) of finite-rank operators $H\to H$.
\end{definition}
\begin{example}
	Let $M$ be a compact space, and choose $K\in C(M\times M)$. Then the integral operator
	\[Af(y)\coloneqq\int_MK(x,y)f(x)\,dx\]
	is compact. By embedding, we may assume $M=[-1,1]^n$. Then we may approximate the integral as a Riemann sum on a tiling by hypercubes of length $1/n$, essentially replacing $K$ by a locally constant function (which then has finite rank).
\end{example}
\begin{lemma}
	Let $A$ be a compact operator on a Hilbert space $H$. If $\{v_n\}_n$ is bounded, then $\{Av_n\}_n$ admits a convegent subsequence.
\end{lemma}
\begin{proof}[Sketch]
	If $A$ has finite rank, there is nothing to do because any sequence in a bounded subset of a finite-dimensional vector space admits a convergent subsequence. In general, choose $\{A_k\}_k$ converging to $k$, and inductively define $v_n^{k+1}$ to be a subsequence of $v_n^k$ (with $v_n^0=v_n$) so that $A_{k+1}v_n^{k+1}$ converges. Then use the diagonal subsequence $v_n^n$ as the required subsequence.
\end{proof}
\begin{theorem}[Hilbert--Schmidt]
	Let $A$ be a self-adjoint compact operator on a Hilbert space $H$. Then there is an orthogonal decomposition
	\[H=\ker A\oplus\widehat{\bigoplus_\lambda}H_\lambda,\]
	where each eigenspace $H_\lambda$ is finite-dimensional, and the set of eigenvalues is finite or tends to zero.
\end{theorem}
\begin{proof}[Sketch]
	If $A$ has finite rank, there is nothing to do. It is essentially enough to prove the theorem for $A^2$ (and then diagonalize everything again against $A$). The key difficulty is exhibiting a single eigenvector for $A$; once we find one eigenvector, we can (roughly speaking) pass to its orthogonal complement and inductively downwards.

	To find our eigenvector, we will actually find the largest eigenvector: we will find a vector in $H$ with eigenvalue $\norm A^2$. Indeed, let $\{A_n\}_n$ be finite-rank operators approaching $A$, and let $v_n$ be an eigenvector of $A_n^2$ with eigenvalue $\norm {A_n}^2$; by passing to a subsequence, we may assume $\{v_n\}$ approaches some $v$, so it follows $A^2v=\lim_nA_nv_n=\norm A^2v$. Now, pass to the orthogonal complement, and our sequence of eigenvalues decreases, and it must tend to zero by compactness of $A$.
\end{proof}

\subsection{The Peter--Weyl Theorem}
\begin{theorem}[Peter--Weyl]
	Let $G$ be a compact group. Then $L^2(G)^{\mathrm{fin}}$ is dense in $L^2(G)$.
\end{theorem}
\begin{proof}[Sketch]
	Let $\{\varphi_n\}_n$ be a continuous, real-valued, symmetric approximate identity. Then $f\mapsto L\varphi_n(f)$ is an integral operator and hence compact, and symmetry and being real-valued implies that they are self-adjoint. Thus, we may diagonalize $L^2(G)$ against these operators. Each nonzero eigenspace is finite-dimensional and $G$-invariant (for right translations) and hence in $L^2(G)^{\mathrm{fin}}$, so $\im L\varphi_n\subseteq L^2(G)^{\mathrm{fin}}$. But for any $f\in L^2(G)$, we have $L\varphi_n(f)\to f$ as $n\to\infty$, so $\sum_n\im L\varphi_n$ is dense in $L^2(G)$!
\end{proof}
\begin{corollary}
	Let $G$ be a compact group. For each irreducible finite-dimensional representation $V$, let $\chi_V\in C(G)$ be the character of $V$. Then $\{\chi_V\}_V$ is an orthonormal basis of the space $L^2(G)^G$ of con\-jugation-invariant functions.
\end{corollary}
\begin{proof}[Sketch]
	Certainly the functions $\{\chi_V\}_V$ are orthnormal by orthogonality of matrix coefficients. To get density, we use the Peter--Weyl theorem: we have a commutative diagram
	% https://q.uiver.app/#q=WzAsNCxbMCwxLCJcXGRpc3BsYXlzdHlsZVxcYmlnb3BsdXNfVihWXFxvdGltZXMgVl4qKSJdLFsxLDEsIkxeMihHKSJdLFswLDAsIlxcZGlzcGxheXN0eWxlXFxiaWdvcGx1c19WKFZcXG90aW1lcyBWXiopXkciXSxbMSwwLCJMXjIoRyleRyJdLFszLDEsIiIsMCx7InN0eWxlIjp7InRhaWwiOnsibmFtZSI6Imhvb2siLCJzaWRlIjoidG9wIn19fV0sWzIsMCwiIiwwLHsic3R5bGUiOnsidGFpbCI6eyJuYW1lIjoiaG9vayIsInNpZGUiOiJ0b3AifX19XSxbMiwzLCJcXHhpIl0sWzAsMSwiXFx4aSJdXQ==&macro_url=https%3A%2F%2Fraw.githubusercontent.com%2FdFoiler%2Fnotes%2Fmaster%2Fnir.tex
	\[\begin{tikzcd}[cramped]
		{\displaystyle\bigoplus_V(V\otimes V^*)^G} & {L^2(G)^G} \\
		{\displaystyle\bigoplus_V(V\otimes V^*)} & {L^2(G)}
		\arrow["\xi", from=1-1, to=1-2]
		\arrow[hook, from=1-1, to=2-1]
		\arrow[hook, from=1-2, to=2-2]
		\arrow["\xi", from=2-1, to=2-2]
	\end{tikzcd}\]
	where each downward inclusion admits a section given by averaging via the integral. Thus, to approximate in the top-right, we push to the bottom-right, approximate by the bottom-left, and then apply the section to get back in the top-left.
\end{proof}

\subsection{Applications}
\begin{proposition}
	Any compact Lie group $G$ admits a faithful finite-dimensional representation.
\end{proposition}
\begin{proof}[Sketch]
	Note that every element of $G$ acts nontrivially on $L^2(G)$ and hence on $L^2(G)^{\mathrm{fin}}$ by density. Thus,
	\[\bigcap_V\ker\rho_V=\{1\},\]
	where the intersection is over all finite-dimensional irreducible reprsentations of $G$. Enumerating the intersection gives a descending sequence of closed Lie subgroups of $G$, which must eventually stabilize because $G$ is a Lie group with finite component group, so some finite intersection equals $\{1\}$.
\end{proof}
\begin{proposition}
	Any compact group is an inverse limit of compact Lie groups.
\end{proposition}
\begin{proof}
	Enumerate the finite-dimensional irreducible representations as $\{V_i\}_i$, and define $W_n\coloneqq\bigoplus_{i\le n}W_i$, and let $K_n$ be the kernel of $W_n$. As in the previous proof, $\bigcap_nK_n=\{1\}$, but now $G/K_n\subseteq\op U(W_n)$ is some compact Lie group, so we are done upon writing $G=\lim_nG/K_n$. 
\end{proof}
\begin{corollary}
	Let $G$ be a compact group. Then $L^2(G)^{\mathrm{fin}}$ is dense in $C(G)$, equipped with $\norm\cdot_\infty$.
\end{corollary}
\begin{proof}
	This follows from the Stone--Weierstrass theorem. The proposition is used to show that $L^2(G)^{\mathrm{fin}}$ separates points.
\end{proof}
\begin{corollary}
	Let $G$ be a compact Lie group, and let $V$ be any faithful finite-dimensional representation of $G$. Then $V\oplus V^*$ $\otimes$-generates $\op{Rep}G$.
\end{corollary}
\begin{proof}
	Define $W\coloneqq V\oplus V^*$; note that $\det W=1$. Because $\op{Rep}G$ is semisimple, it is enough to realize all irreducible representations of $G$, so it is enough to realize $L^2(G)^{\mathrm{fin}}$. In fact, we claim that the matrix coefficient map $\op{Sym}(W\otimes W^*)\to L^2(G)^{\mathrm{fin}}$ is surjective: the generated subalgebra $A$ is left-invariant, separates points (because $W$ is faithful), and it is stable under complex conjugation because any $g\in G$ has $\rho_W(g)^\dagger=\rho_W(g)^{-1}=\rho_{\land^{\dim W-1}W}(g)$. Thus, $A$ is dense in $C(G)$, so we find $A=L^2(G)^{\mathrm{fin}}$ by computing the multiplicity spaces $\op{Hom}_G(V,A)$ as before to be $V^*$.
\end{proof}

\section{Real Forms}

\subsection{Counting with Cocycles}
\begin{proposition}
	Fix a Galois field extension $L/K$ and a Lie algebra $\mf g$ over $K$. Then $\mathrm H^1(\op{Gal}(L/K);\op{Aut}\mf g_L)$ classifies Lie algebras $\mf g'$ over $K$ such that $\mf g'_L\cong\mf g_L$.
\end{proposition}
\begin{proof}[Sketch]
	Given a $1$-cocycle $c\colon\op{Gal}(L/K)\to\op{Aut}\mf g_L$, recover $\mf g'$ by taking Galois-invariants. Conversely, given $\mf g'$, define the $1$-cocycle by sending $\sigma\in\op{Gal}(L/K)$ to the automorphism given by moving around the following square.
	% https://q.uiver.app/#q=WzAsNCxbMCwwLCJcXG1mIGcnX0wiXSxbMSwwLCJcXG1mIGdfTCJdLFswLDEsIlxcbWYgZydfTCJdLFsxLDEsIlxcbWYgZ19MIl0sWzAsMiwiXFxzaWdtYSIsMl0sWzEsMywiXFxzaWdtYSJdLFswLDEsIlxcc2ltIl0sWzIsMywiXFxzaW0iXV0=&macro_url=https%3A%2F%2Fraw.githubusercontent.com%2FdFoiler%2Fnotes%2Fmaster%2Fnir.tex
	\[\begin{tikzcd}[cramped]
		{\mf g'_L} & {\mf g_L} \\
		{\mf g'_L} & {\mf g_L}
		\arrow["\sim", from=1-1, to=1-2]
		\arrow["\sigma"', from=1-1, to=2-1]
		\arrow["\sigma", from=1-2, to=2-2]
		\arrow["\sim", from=2-1, to=2-2]
	\end{tikzcd}\]
	Modifying the choice of isomorphism adjusts the $1$-cocycle by a corresponding $1$-coboundary.
\end{proof}
\begin{corollary}
	Fix a semisimple complex Lie algebra $\mf g$. Then real forms of $\mf g$, denoted $\mf g_{(s)}$, are classified by $s\in\op{Aut}\mf g$ satisfying $s\overline s$, up to the conjugation action $a\cdot s=as\overline a^{-1}$.
\end{corollary}
\begin{proof}
	Apply the proposition to the split real form of $\mf g$. In particular, define $s$ to be the image of the nontrivial element of $\op{Gal}(\CC/\RR)$ along the $1$-cocycle.
\end{proof}
\begin{remark}
	By construction, $\mf g_{(s)}=\{X\in\mf g_\CC:X=\overline{sX}\}$.
\end{remark}
\begin{definition}[inner class]
	Two real forms $\mf g_{(s)}$ and $\mf g_{(s')}$ are in the same \textit{inner class} if and only if $s$ and $s'$ are in the same connected component of $\op{Aut}\mf g$ (up to conjugation).
\end{definition}
\begin{definition}[quasisplit]
	The real form $\mf g_{(s)}$ is \textit{quasisplit} if and only if the component of $s$ in $G_{\mathrm{ad}}$ is trivial (up to conjugation).
\end{definition}
\begin{definition}[compact form]
	The \textit{compact real form} $\mf g^c$ is defined $\mf g_{(\tau)}$, where $\tau$ is the involution defined on the Serre presentation by
	\[\begin{cases}
		\tau(h_i)=h_i, \\
		\tau(e_i)=-f_i, \\
		\tau(f_i)=-e_i.
	\end{cases}\]
\end{definition}
\begin{remark}
	The Killing form on $\mf g^c$ is negative-definite: one can divide this calculation into $\mf h\cap\mf g^c$ (handled by direct calculatoin) and the pieces $(\mf g_\alpha\oplus\mf g_{-\alpha})\cap\mf g^c$ (handled by using the Weyl group to reduce to the case of simple roots and then to $\mf{sl}_2$). Thus, the corrsponding subgroup $G_{\mathrm{ad}}^c\subseteq G_{\mathrm{ad}}$ is compact.
\end{remark}
\begin{remark}
	A similar calculation on the root decomposition shows that the Killing form on the split real form has signature $(\op{rank}\mf g+\#\Phi^+,\#\Phi^+)$.
\end{remark}
\begin{remark}
	The image of $\tau$ in $\op{Aut}D$ is $-w_0$, where $w_0$ is the long Weyl element. Indeed, $\tau$ acts by $(-1)$ on the root system.
\end{remark}
\begin{proposition}
	Fix a semisimple complex Lie algebra $\mf g$. Then real forms $\mf g_\theta$ are classified by conjugacy classes of involutions $\theta$ commuting with $X\mapsto\overline{\tau X}$.
\end{proposition}
\begin{proof}[Sketch]
	Set $\omega\coloneqq\overline\tau$ for brevity. Then other antilinear involutions take the form $\omega\circ g$ for some $g\in\op{Aut}\mf g$. By replacing $g$ with $\theta\coloneqq g\left|g\right|$ (note $\theta=\left|g\right|^{-1/2}g\omega\left(\left|g\right|^{1/2}\right)$ also), we may assume $\theta^2=1$. Similarly judiciously taking square roots shows that the conjugacy relations match up as well.
\end{proof}
\begin{definition}
	Let $\mf g_\theta$ be a real form of $\mf g$. Then we define $G_{\mathrm{ad},\theta}\coloneqq G_{\mathrm{ad}}^{\omega_\theta}$, where $\omega_\theta\coloneqq\overline\tau\circ\theta$ is the corresponding antilinear involution.
\end{definition}

\subsection{The Vogan Game}
\begin{remark}
	Given $\theta$, let $\mf g=\mf k\oplus\mf p$ be the eigenspace decomposition, where $\mf k=\mf g^{\theta=1}$ and $\mf p=\mf g^{\theta=-1}$. Then by construction, $\mf g_\theta=\mf k^c\oplus i\mf p^c$. This is an orthogonal decomposition for the Killing form.
\end{remark}
\begin{example}
	The signature of the Killing form on $\mf g_\theta$ is $(\dim\mf p,\dim\mf k)$. It follows that the split real form has $\dim\mf k=\#\Phi^+$.
\end{example}
\begin{definition}[Vogan diagram]
	A \textit{Vogan diagram} is a Dynkin diagram with the following decorations: each vertex is either
	\begin{itemize}
		\item connected with an auxiliary edge corresponding to an involution of the Dynkin diagram, or
		\item colored black or white.
	\end{itemize}
\end{definition}
\begin{proposition}
	Each real form $\mf g_\theta$ produces a Vogan diagram, and each Vogan diagram comes from a real form.
\end{proposition}
\begin{proof}[Sketch]
	By choosing a Cartan subalgebra $\mf h_+$ of $\mf k^c$ first and then extending it maximally by $\mf h_-$ to $i\mf p^c$, one can find a $\theta$-stable Cartan subalgebra of $\mf g^c$. It turns out that $\mf h_-$ contains no coroots, so a generic $t\in\mf h_+$ defines a polarization. Thus, $\theta$ now acts on the root decomposition preserving the positive roots and thus the simple roots. Thus, it induces an automorphism of the Dynkin diagram; by appropriate normalization, if $\theta(\alpha_i)\ne\alpha_i$, we may define $e_{\theta(i)}\coloneqq\theta(e_i)$, and this corresponds to connected roots. Otherwise, $\theta(i)=i$, so $\theta(e_i)=\pm e_i$. The $+$ case corresponds to a white color, and the $-$ case corresponds to a black color. One can read the Vogam diagram backwards to recover $\theta$.
\end{proof}
\begin{example}
	If $\mf g_\theta$ is in the compact inner class, then $\theta$ must act by identity on the Dynkin diagram, so $\mf h\subseteq\mf k$, so $\op{rank}\mf k=\op{rank}\mf g$.
\end{example}
\begin{remark}[Vogan game]
	Suppose that the vertex $\alpha_i$ is black. Then conjugating $\theta$ by $s_i$ adjusts the sign of $e_j$ from $\alpha_j(\theta)$ to $\alpha_j(s_i\theta)$, which is $s_i\alpha_j(\theta)$, which is
	\[\alpha_j(\theta)\alpha_i(\theta)^{-a_{ij}}.\]
	Thus, exactly the vertices $\alpha_j$ adjacent to $\alpha_i$ with $a_{ij}$ (namely, pointing in the correct direction with odd multiplicity) change color.
\end{remark}
\begin{remark}
	If all vertices of a Vogan diagram (in the compact inner class) are white except for $\alpha_i$, and $\omega_i$ is miniscule, then $\mf k$ is the sum of $\mf{gl}_1$ and the white vertices. Indeed, $\mf k$ has all $\mf h$, and it features the root spaces $\mf g_\alpha$ which feature $\alpha_i$ an even number of times. This means that $(\alpha,\omega_i)=0$ because $\omega_i$ is miniscule.
\end{remark}
\begin{remark}
	If all vertices of a Vogan diagram (in the compact inner class) are white except for $\alpha_i$, and $\omega_i$ is maximal, then $\mf k$ is the sum of $\mf{sl}_{2}$ and the white vertices. Indeed, one argues as in the previous remark, using maximality to note that the only root $\alpha$ featuring $\alpha_i$ at least twice is $\omega_i$.
\end{remark}
\begin{remark}
	If all vertices of the Vogan diagram are white, then $\mf k$ is given by ``folding'' the Dynkin diagram in half (and adjusting lengths appropriately). Indeed, the new root system features the simple roots $\{\alpha_i:\theta(i)=i\}$ and $\{\alpha_i+\alpha_{\theta(i)}:\theta(i)\ne i\}$.
\end{remark}

\subsection{Polar Decomposition}
\begin{definition}
	Let $\mf g_\theta$ be a real form of $\mf g$. Then define $K^c_\theta$ to be the fixed points of $\omega_\theta\coloneqq\omega\circ\theta$ in $G_{\mathrm{ad},\theta}$.
\end{definition}
\begin{remark}
	One can expand out the definition to show that $\omega(g)=g^{-\dagger}$, where the implicity Hermitian form is the unique Hermitian extension of the Killing form restricted to $\mf g^c$. Thus, $K^c_\theta=(G_{\mathrm{ad}}^c)^{\omega_\theta}$ is subgroup of $G_{\mathrm{ad},\theta}$ acting by unitary operators.
\end{remark}
\begin{definition}
	Let $\mf g_\theta$ be a real form of $\mf g$. Then define $P_\theta\coloneqq\exp(\mf p_\theta)$, where $\mf p_\theta\coloneqq i\mf p^c$.
\end{definition}
\begin{remark}
	One can show that $X\in\mf g$ acts by a Hermitian operator on $\mf g$ if and only if $\omega(X)=-X$. Thus, $X\in\mf g_\theta$ acts by a Hermitian operator if and only if $X\in i\mf p^c$. Because the exponential is a diffeomorphism from Hermitian operators to positive Hermitian operators, we conclude that $P_\theta\subseteq G_{\mathrm{ad},\theta}$ is the submanifold of positive Hermitian operatots.
\end{remark}
\begin{theorem}[Polar decomposition]
	The multiplication map $K^c_\theta\times P_\theta\to G_{\mathrm{ad},\theta}$ is a diffeomorphism.
\end{theorem}
\begin{proof}[Sketch]
	We construct an explicit smooth inverse. The idea is to do a polar decomposition in $\op{GL}(\mf g)$. Namely, the operator $\op{Ad}_g$ is uniquely factored as $U_gR_g$, where $U_g\in\op{GL}(\mf g)$ is unitary and $R_g\in\op{GL}(\mf g)$ is positive Hermitian; we would like to show that $g\mapsto(U_g,R_g)$ is our inverse. Well, $R_g\coloneqq(\mathrm{Ad}_g^\dagger\mathrm{Ad}_g)^{1/2}$, and an investigation of the action on the root decomposition reveals $R_g\in\op{Aut}(\mf g)$ and then $R_g\in P_\theta$. Thus, $U_g\in G_{\mathrm{ad},\theta}$ and so $U_g\in K^c_\theta$.
\end{proof}

\section{Compact Groups}

\subsection{The Fundamental Group}
\begin{lemma}
	The group $G_{\mathrm{ad}}$ has trivial center.
\end{lemma}
\begin{proof}
	Suppose $z\in G_{\mathrm{ad}}$ is central. Then $\op{Ad}_z\colon\mf g\to\mf g$ is a scalar, so $z\in\op{GL}(\mf g)$ is a scalar matrix. However, this means $z[X,Y]=[zX,zY]$ for all $X,Y\in\mf g$, so $z=1$ follows.
\end{proof}
\begin{proposition}
	Let $G$ be a complex Lie group with semisimple Lie algebra $\mf g$ and center $Z$. For each dominant integral weight $\lambda$, let $\chi_\lambda\colon Z\to\CC^\times$ be the central character of $L_\lambda$. Then $\chi_\bullet$ extends uniquely to an injective homomorphism
	\[P/Q\to\op{Hom}(Z,\CC^\times).\]
\end{proposition}
\begin{proof}[Sketch]
	To show that $\chi_\bullet$ extends to a homomorphism $P\to\op{Hom}(Z,\CC^\times)$, note that $L_{\lambda+\mu}\subseteq L_\lambda\otimes L_\mu$, so $\chi_{\lambda+\mu}=\chi_\lambda\chi_\mu$ for dominant integral $\lambda$ and $\mu$. It remains to compute the kernel.
	\begin{itemize}
		\item If $\chi_\lambda=1$, then $L_\lambda$ descends to a representation of $G/Z$, which is $G_{\mathrm{ad}}$. But $\mf g$ is a faithful unimodular representation of $G_{\mathrm{ad}}^c$, so $\mf g$ $\otimes$-generates $\op{Rep}G_{\mathrm{ad}}^c$. Thus, $L_\lambda$ may only feature weights from $Q$, so $\lambda\in Q$.
		\item To show $Q\subseteq\ker\chi_\bullet$, it's enough to show that, for any root $\alpha$, we have $\chi_\lambda\chi_\theta=\chi_{\lambda+\alpha}$ for large enough dominant integral weights $\lambda$. In fact, $L_{\lambda+\alpha}\subseteq L_\lambda\otimes\mf g$ for $\lambda$ large enough: by unwinding highest weight theory and the Serre presentation, we see that the space of maps $L_{\lambda+\alpha}\to\mf g\otimes L_\lambda$ is given by
		\[\mf g[\alpha]_\lambda=\left\{v\in\mf g[\alpha]:f_i^{(\lambda,\alpha_i^\lor)}v=0\text{ for all }i\right\},\]
		but the condition is automatically true if $\lambda$ is large.
		\qedhere
	\end{itemize}
\end{proof}

\newpage
\appendix

\section{Type \texorpdfstring{$A_n$}{An}}
\begin{itemize}
	\item Lie algebra
	\begin{itemize}
		\item $\mf g=\mf{sl}(n+1)$.
		\item $\mf h=\{\op{diag}(a_1,\ldots,a_{n+1}):a_1+\ldots+a_{n+1}=0\}$.
		\item $\mf h^*$ is spanned by the projections $e_i\colon\op{diag}(a_1,\ldots,a_{n+1})\mapsto a_i$, where $e_1+\cdots+e_{n+1}=0$.
		\item Killing form $K(X,Y)=2(n+1)\tr XY$.
	\end{itemize}
	\item Root system
	\[\begin{array}{c|cccccc}
		\text{root} & \text{root space} & \text{coroot} \\\hline
		e_i-e_j & E_{ij} & E_{ii}-E_{jj}\\
	\end{array}\]
	\item Lattices
	\begin{itemize}
		\item Weight lattice $P=\{(\lambda_1,\ldots,\lambda_{n+1})\in\RR^{n+1}/(1,\ldots,1)\RR:\lambda_i-\lambda_{i+1}\in\ZZ\text{ for all }i\}$.
		\item Dominant weights $P_+=\{(\lambda_1,\ldots,\lambda_n)\in P:\lambda_1\ge\cdots\ge\lambda_n\}$.
		\item Root lattice $Q=\{(\lambda_1,\ldots,\lambda_{n+1})\in\ZZ^{n+1}:\lambda_1+\cdots+\lambda_{n+1}=0\}$.
		\item $P/Q\cong\ZZ/(n+1)\ZZ$ by $(\lambda_1,\ldots,\lambda_{n+1})\mapsto\sum_i\lambda_i$.
	\end{itemize}
	\item Polarization
	\begin{itemize}
		\item The positive roots are $\{e_i-e_j:i<j\}$.
		\item The simple roots are $\{e_i-e_{i+1}:i\}$. Set $\alpha_i\coloneqq e_i-e_{i+1}$.
		\item Dynkin diagram:
		% https://q.uiver.app/#q=WzAsNCxbMCwwLCJcXGFscGhhXzEiXSxbMSwwLCJcXGFscGhhXzIiXSxbMiwwLCJcXGNkb3RzIl0sWzMsMCwiXFxhbHBoYV9uIl0sWzAsMSwiIiwwLHsic3R5bGUiOnsiaGVhZCI6eyJuYW1lIjoibm9uZSJ9fX1dLFsxLDIsIiIsMCx7InN0eWxlIjp7ImhlYWQiOnsibmFtZSI6Im5vbmUifX19XSxbMiwzLCIiLDAseyJzdHlsZSI6eyJoZWFkIjp7Im5hbWUiOiJub25lIn19fV1d&macro_url=https%3A%2F%2Fraw.githubusercontent.com%2FdFoiler%2Fnotes%2Fmaster%2Fnir.tex
		\[\begin{tikzcd}[cramped]
			{\alpha_1} & {\alpha_2} & \cdots & {\alpha_n}
			\arrow[no head, from=1-1, to=1-2]
			\arrow[no head, from=1-2, to=1-3]
			\arrow[no head, from=1-3, to=1-4]
		\end{tikzcd}\]
		\item The fundamental weights are $\{\omega_i\}_i$ where $\omega_i\coloneqq({1,\ldots,1},0,\ldots,0)$ has $i$ ones.
		\item Maximal root $\theta=e_{n+1}-e_n$.
	\end{itemize}
	\item Weyl group
	\begin{itemize}
		\item Weyl group is $W=S_{n+1}$ acting by permutations of coordinates.
		\item The simple reflection $s_i$ is $(i,i+1)$.
		\item $w_0=(n+1,n,\ldots,2,1)$.
	\end{itemize}
	\item Representation theory
	\begin{itemize}
		\item $\rho=(n/2,n/2-1,\ldots,1-n/2,-n/2)$ and $\rho^\lor=(n/2,n/2-1,\ldots,1-n/2,-n/2)$.
		\item The miniscule weights are $\{0\}\sqcup\{\omega_i:i\}$.
		\item Fundamental representations are $L_{\omega_i}=\land^i\CC^{n+1}$.
		\item Exponents are $\{1,2,\ldots,n\}$.
	\end{itemize}
	\item Group theory
	\begin{itemize}
		\item $G_{\mathrm{sc}}=\op{SL}_{n+1}(\CC)$.
		\item $G_{\mathrm{ad}}=\op{PSL}_{n+1}(\CC)$.
	\end{itemize}
\end{itemize}

\newpage
\section{Type \texorpdfstring{$B_n$}{ Bn}}
\begin{itemize}
	\item Lie algebra
	\begin{itemize}
		\item $\mf g=\mf{so}(2n+1)$, preserving
		\[B=\begin{bmatrix}
			& 1_n \\ 1_n \\ && 1
		\end{bmatrix}.\]
		\item $\mf h=\{\op{diag}(a_1,\ldots,a_n,-a_1,\ldots,-a_n,0):a_1,\ldots,a_n\}$.
		\item $\mf h^*$ is spanned by the projections $e_i\colon\op{diag}(a_1,\ldots,a_n,-a_1,\ldots,-a_n,0)\mapsto a_i$.
		\item Killing form $K(X,Y)=2(n+1)\tr XY$.\todo{}
	\end{itemize}
	\item Root system
	\[\begin{array}{c|cccccc}
		\text{root} & \text{length} & \text{root space} & \text{coroot}\\\hline
		e_i-e_j & \text{long} & E_{ij}-E_{i+n,j+n} & H_i-H_j\\
		e_i+e_j & \text{long} & E_{i,j+n}-E_{j,i+n} & H_i+H_j\\
		-e_i-e_j & \text{long} & E_{i+n,j}-E_{j+n,i} & -H_i-H_j \\
		e_i & \text{short} & E_{i,2n+1}-E_{2n+1,i+n} & 2H_i \\
		-e_i & \text{short} & E_{n+i,2n+1}-E_{2n+1,i} & -2H_i
	\end{array}\]
	Here, $H_i=E_{ii}-E_{i+n,i+n}$.
	\item Lattices
	\begin{itemize}
		\item Weight lattice $P=\{(\lambda_1,\ldots,\lambda_{n})\in(\frac12\ZZ)^n:\lambda_i-\lambda_{i+1}\in\ZZ\text{ for all }i\}$.
		\item Dominant weights $P_+=\{(\lambda_1,\ldots,\lambda_n)\in P:\lambda_1\ge\cdots\ge\lambda_n\ge0\}$.
		\item Root lattice $Q=\ZZ^n$.
		\item $P/Q\cong\ZZ/2\ZZ$ by $(\lambda_1,\ldots,\lambda_n)\mapsto\lambda_1$.
	\end{itemize}
	\item Polarization
	\begin{itemize}
		\item The positive roots are $\{e_i-e_j:i<j\}\cup\{e_i+e_j:i<j\}\cup\{e_i:i\}$.
		\item The simple roots are $\{e_i-e_{i+1}:i\}\cup\{e_n\}$. Set $\alpha_i\coloneqq e_i-e_{i+1}$ and $\alpha_n\coloneqq e_n$.
		\item Dynkin diagram:
		\[\begin{tikzcd}[cramped]
			{\alpha_1} & {\alpha_2} & \cdots & {\alpha_{n-1}} & {\alpha_n}
			\arrow[no head, from=1-1, to=1-2]
			\arrow[no head, from=1-2, to=1-3]
			\arrow[no head, from=1-3, to=1-4]
			\arrow[Leftarrow, nfold, from=1-5, to=1-4]
		\end{tikzcd}\]
		\item The fundamental weights are $\{\omega_i\}_i$ where $\omega_i\coloneqq({1,\ldots,1},0,\ldots,0)$ has $i$ ones for $i<n$ and $\omega_n\coloneqq\frac12(1,\ldots,1)$.
		\item Maximal root $\theta=e_1+e_2$.
	\end{itemize}
	\item Combinatorics
	\begin{itemize}
		\item Weyl group is $W=S_{n}\ltimes\{\pm1\}^n$ acting by permutations of coordinates and sign changes.
		\item The simple reflection $s_i$ is $(i,i+1)$ for $i<n$, and $s_n$ is $(+1,\ldots,+1,-1)$.
		\item $w_0=(-1,\ldots,-1)$.
	\end{itemize}
	\item Representation theory
	\begin{itemize}
		\item $\rho=(n-1/2,n-3/2,\ldots,3/2,1/2)$ and $\rho^\lor=(n,n-1,\ldots,2,1)$.
		\item The miniscule weights are $\{0\}\sqcup\{\omega_1\}$.
		\item Fundamental representations are $L_{\omega_i}=\land^i\CC^{2n+1}$ for $i<n$ and the standard Clifford module for $i=n$.
		\item Exponents are $\{1,3,\ldots,2n-3,2n-1\}$.
	\end{itemize}
	\item Group theory
	\begin{itemize}
		\item $G_{\mathrm{sc}}=\op{Spin}_{2n+1}(\CC)$.
		\item $G_{\mathrm{ad}}=\op{SO}_{2n+1}(\CC)$.
	\end{itemize}
\end{itemize}

\newpage
\section{Type \texorpdfstring{$C_n$}{ Cn}}
\begin{itemize}
	\item Lie algebra
	\begin{itemize}
		\item $\mf g=\mf{sp}(2n)$, preserving
		\[B=\begin{bmatrix}
			& -1_n \\ 1_n
		\end{bmatrix}.\]
		\item $\mf h=\{\op{diag}(a_1,\ldots,a_n,-a_1,\ldots,-a_n):a_1,\ldots,a_n\}$.
		\item $\mf h^*$ is spanned by the projections $e_i\colon\op{diag}(a_1,\ldots,a_n,-a_1,\ldots,-a_n)\mapsto a_i$.
		\item Killing form $K(X,Y)=2(n+1)\tr XY$.\todo{}
	\end{itemize}
	\item Root system
	\[\begin{array}{c|cccccc}
		\text{root} & \text{length} & \text{root space} & \text{coroot} \\\hline
		e_i-e_j & \text{short} & E_{ij}-E_{i+n,j+n} & H_i-H_j \\
		e_i+e_j & \text{short} & E_{i,j+n}-E_{j,i+n} & H_i+H_j \\
		-e_i-e_j & \text{short} & E_{i+n,j}-E_{j+n,i} & -H_i-H_j \\
		2e_i & \text{long} & E_{i,i+n} & H_i \\
		-2e_i & \text{long} & E_{i+n,i} & -H_i
	\end{array}\]
	Here, $H_i=E_{ii}-E_{i+n,i+n}$.
	\item Lattices
	\begin{itemize}
		\item Weight lattice $P=\ZZ^n$.
		\item Dominant weights $P_+=\{(\lambda_1,\ldots,\lambda_n)\in P:\lambda_1\ge\cdots\ge\lambda_n\ge0\}$.
		\item Root lattice $Q=\{(\lambda_1,\ldots,\lambda_n)\in\ZZ^n:\lambda_1+\cdots+\lambda_n\equiv0\pmod2\}$.
		\item $P/Q\cong\ZZ/2\ZZ$ by $(\lambda_1,\ldots,\lambda_n)\mapsto\sum_i\lambda_i$.
	\end{itemize}
	\item Polarization
	\begin{itemize}
		\item The positive roots are $\{e_i-e_j:i<j\}\cup\{e_i+e_j:i<j\}\cup\{2e_i:i\}$.
		\item The simple roots are $\{e_i-e_{i+1}:i\}\cup\{2e_n\}$. Set $\alpha_i\coloneqq e_i-e_{i+1}$ and $\alpha_n\coloneqq2e_n$.
		\item Dynkin diagram:
		% https://q.uiver.app/#q=WzAsNSxbMCwwLCJcXGFscGhhXzEiXSxbMSwwLCJcXGFscGhhXzIiXSxbMiwwLCJcXGNkb3RzIl0sWzMsMCwiXFxhbHBoYV97bi0xfSJdLFs0LDAsIlxcYWxwaGFfbiJdLFswLDEsIiIsMCx7InN0eWxlIjp7ImhlYWQiOnsibmFtZSI6Im5vbmUifX19XSxbMSwyLCIiLDAseyJzdHlsZSI6eyJoZWFkIjp7Im5hbWUiOiJub25lIn19fV0sWzIsMywiIiwwLHsic3R5bGUiOnsiaGVhZCI6eyJuYW1lIjoibm9uZSJ9fX1dLFs0LDMsIiIsMix7ImxldmVsIjoyfV1d&macro_url=https%3A%2F%2Fraw.githubusercontent.com%2FdFoiler%2Fnotes%2Fmaster%2Fnir.tex
		\[\begin{tikzcd}[cramped]
			{\alpha_1} & {\alpha_2} & \cdots & {\alpha_{n-1}} & {\alpha_n}
			\arrow[no head, from=1-1, to=1-2]
			\arrow[no head, from=1-2, to=1-3]
			\arrow[no head, from=1-3, to=1-4]
			\arrow[Rightarrow, nfold, from=1-5, to=1-4]
		\end{tikzcd}\]
		\item The fundamental weights are $\{\omega_i\}_i$ where $\omega_i\coloneqq({1,\ldots,1},0,\ldots,0)$ has $i$ ones.
		\item Maximal root $\theta=2e_1$.
	\end{itemize}
	\item Combinatorics
	\begin{itemize}
		\item Weyl group is $W=S_{n}\ltimes\{\pm1\}^n$ acting by permutations of coordinates and sign changes.
		\item The simple reflection $s_i$ is $(i,i+1)$ for $i<n$, and $s_n$ is $(+1,\ldots,+1,-1)$.
		\item $w_0=(-1,\ldots,-1)$.
	\end{itemize}
	\item Representation theory
	\begin{itemize}
		\item $\rho=(n,n-1,\ldots,2,1)$ and $\rho^\lor=(n-1/2,n-3/2,\ldots,3/2,1/2)$.
		\item The miniscule weights are $\{0\}\sqcup\{\omega_1\}$.
		\item Fundamental representations are $L_{\omega_i}=\ker\left(\iota_B\colon\land^i\CC^{2n}\to\land^{i-2}\CC^{2n}\right)$, where $\iota_B$ is contraction:
		\[\iota_B(v_1\land\cdots\land v_i)\coloneqq\sum_{r<s}(-1)^{r+s+1}B(v_r\land v_s)(v_1\land\cdots\land\widehat v_r\land\cdots\land\widehat v_s\land\cdots\land v_i).\]
		\item Exponents are $\{1,3,\ldots,2n-3,2n-1\}$.
	\end{itemize}
	\item Group theory
	\begin{itemize}
		\item $G_{\mathrm{sc}}=\op{Sp}_{2n}(\CC)$.
		\item $G_{\mathrm{ad}}=\op{Sp}_{2n}(\CC)/\{\pm1\}$.
	\end{itemize}
\end{itemize}

\newpage
\section{Type \texorpdfstring{$D_n$}{ Dn}}
\begin{itemize}
	\item Lie algebra
	\begin{itemize}
		\item $\mf g=\mf{so}(2n)$, preserving
		\[B=\begin{bmatrix}
			& 1_n \\ 1_n
		\end{bmatrix}.\]
		\item $\mf h=\{\op{diag}(a_1,\ldots,a_n,-a_1,\ldots,-a_n):a_1,\ldots,a_n\}$.
		\item $\mf h^*$ is spanned by the projections $e_i\colon\op{diag}(a_1,\ldots,a_n,-a_1,\ldots,-a_n)\mapsto a_i$.
		\item Killing form $K(X,Y)=2(n+1)\tr XY$.\todo{}
	\end{itemize}
	\item Root system
	\[\begin{array}{c|cccccc}
		\text{root} & \text{root space} & \text{coroot} \\\hline
		e_i-e_j & E_{ij}-E_{j+n,i+n} & H_i-H_j \\
		e_i+e_j & E_{i,j+n}-E_{j,i+n} & H_i+H_j \\
		-e_i-e_j & E_{i+n,j}-E_{j+n,i} & -H_i-H_j \\
	\end{array}\]
	Here, $H_i=E_{ii}-E_{i+n,i+n}$.
	\item Lattices
	\begin{itemize}
		\item Weight lattice $P=\{(\lambda_1,\ldots,\lambda_{n})\in(\frac12\ZZ)^n:\lambda_i-\lambda_{i+1}\in\ZZ\text{ for all }i\}$.
		\item Dominant weights $P_+=\{(\lambda_1,\ldots,\lambda_n)\in P:\lambda_1\ge\cdots\ge\lambda_n,\lambda_{n-1}+\lambda_n\ge0\}$.
		\item Root lattice $P=\{(\lambda_1,\ldots,\lambda_{n})\in\ZZ^n:\lambda_1+\cdots+\lambda_n\equiv0\pmod2\}$.
		\item $P/Q\cong(\ZZ/2\ZZ)^2$ if $n$ even and $P/Q\cong\ZZ/4\ZZ$ if $n$ odd. It is generated by $(1,0,\ldots,0)$ and $\frac12(1,\ldots,1)$.
	\end{itemize}
	\item Polarization
	\begin{itemize}
		\item The positive roots are $\{e_i-e_j:i<j\}\cup\{e_i+e_j:i<j\}$.
		\item The simple roots are $\{e_i-e_{i+1}:i\}\cup\{e_{n-1}+e_n\}$. Set $\alpha_i\coloneqq e_i-e_{i+1}$ and $\alpha_n\coloneqq e_{n-1}+e_n$.
		\item Dynkin diagram:
		% https://q.uiver.app/#q=WzAsNixbMCwxLCJcXGFscGhhXzEiXSxbMSwxLCJcXGFscGhhXzIiXSxbMiwxLCJcXGNkb3RzIl0sWzMsMSwiXFxhbHBoYV97bi0yfSJdLFs0LDAsIlxcYWxwaGFfe24tMX0iXSxbNCwyLCJcXGFscGhhX24iXSxbMCwxLCIiLDAseyJzdHlsZSI6eyJoZWFkIjp7Im5hbWUiOiJub25lIn19fV0sWzEsMiwiIiwwLHsic3R5bGUiOnsiaGVhZCI6eyJuYW1lIjoibm9uZSJ9fX1dLFsyLDMsIiIsMCx7InN0eWxlIjp7ImhlYWQiOnsibmFtZSI6Im5vbmUifX19XSxbMyw0LCIiLDAseyJzdHlsZSI6eyJoZWFkIjp7Im5hbWUiOiJub25lIn19fV0sWzMsNSwiIiwwLHsic3R5bGUiOnsiaGVhZCI6eyJuYW1lIjoibm9uZSJ9fX1dXQ==&macro_url=https%3A%2F%2Fraw.githubusercontent.com%2FdFoiler%2Fnotes%2Fmaster%2Fnir.tex
		\[\begin{tikzcd}[cramped, row sep=tiny]
			&&&& {\alpha_{n-1}} \\
			{\alpha_1} & {\alpha_2} & \cdots & {\alpha_{n-2}} \\
			&&&& {\alpha_n}
			\arrow[no head, from=2-1, to=2-2]
			\arrow[no head, from=2-2, to=2-3]
			\arrow[no head, from=2-3, to=2-4]
			\arrow[no head, from=2-4, to=1-5]
			\arrow[no head, from=2-4, to=3-5]
		\end{tikzcd}\]
		\item The fundamental weights are $\{\omega_i\}_i$ where $\omega_i\coloneqq({1,\ldots,1},0,\ldots,0)$ has $i$ ones for $i<n-1$ and $\omega_{n-1}\coloneqq\frac12(1,\ldots,1,-1)$ and $\omega_n\coloneqq\frac12(1,\ldots,1)$.
		\item Maximal root $\theta=e_1+e_2$.
	\end{itemize}
	\item Combinatorics
	\begin{itemize}
		\item Weyl group is $W=S_{n}\ltimes\{\pm1\}_+^n$ acting by permutations of coordinates and even sign changes.
		\item The simple reflection $s_i$ is $(i,i+1)$ for $i<n$, and $s_n$ is $(+1,\ldots,+1,-1,-1)$.
		\item If $n$ even, $w_0=(-1,\ldots,-1)$ acts by $(-1)$ on $P/Q$. If $n$ odd, $w_0=(-1,\ldots,-1,+1)$ does not.
	\end{itemize}
	\item Representation theory
	\begin{itemize}
		\item $\rho=(n-1,n-2,\ldots,1,0)$ and $\rho^\lor=(n-1,n-2,\ldots,1,0)$.
		\item The miniscule weights are $\{0\}\sqcup\{\omega_1,\omega_{n-1},\omega_n\}$.
		\item Fundamental representations are $L_{\omega_i}=\land^i\CC^{2n+1}$ for $i<n$ and the degree pieces of the standard Clifford module.
		\item Exponents are $\{1,3,\ldots,2n-3,2n-1\}\sqcup\{n\}$.
	\end{itemize}
	\item Group theory
	\begin{itemize}
		\item $G_{\mathrm{sc}}=\op{Spin}_{2n}(\CC)$.
		\item $G_{\mathrm{ad}}=\op{SO}_{2n}(\CC)/\{\pm1\}$.
	\end{itemize}
\end{itemize}

\end{document}